\section{CASPT2: the class table and the three contraction paths}
\label{sec:caspt2-tutorial}

The perturbation layer reports its energy class by class:
\texttt{V0\_ijab}, \texttt{Vm1\_iab}, \texttt{Vm2\_ab},
\texttt{V1\_ija}, \texttt{V2\_ij}, \texttt{V0\_ia},
\texttt{Vm1\_a}, and \texttt{V1\_i}.  Reading that table is the
difference between knowing that a number is close to a reference and
knowing which part of the wavefunction produced it.  This chapter
builds a small reference, runs the three contraction routes, and shows
how to compare them.

\subsection{Build a reference}

The reference chain is RHF, CASSCF, state-specific CASCI, then
canonicalization.  The canonicalization is not cosmetic: the CASPT2
denominators are defined in the semicanonical basis, and the
active-natural-orbital rotation is part of the definition.

\begin{lstlisting}[language=Python]
import numpy as np
import jax
jax.config.update('jax_enable_x64', True)
import jax.numpy as jnp

from pyscf import gto, scf, mcscf
from diffmrc.core.casci import cs_casci
from diffmrc.core.nevpt2.impl import canonicalize_jax, cas_dets_for2

ncore, nact, nele = 4, 2, 2
nocc = ncore + nact

mol = gto.M(atom='O 0 0 0; H 0.7557 0 0.5880; H -0.7557 0 0.5880',
            basis='6-31g', unit='Angstrom', verbose=0)
mf = scf.RHF(mol).run(verbose=0)
mc = mcscf.CASSCF(mf, nact, nele)
mc.kernel()

C = jnp.asarray(mc.mo_coeff)
h1e = jnp.asarray(mol.intor('int1e_kin') + mol.intor('int1e_nuc'))
eri = jnp.asarray(mol.intor('int2e', aosym='s1'))

def mo_ints(C, h1e, eri):
    h1 = jnp.einsum('ai,ab,bj->ij', C, h1e, C)
    z = jnp.einsum('pi,pqrs->iqrs', C, eri)
    z = jnp.einsum('qj,iqrs->ijrs', C, z)
    z = jnp.einsum('rk,ijrs->ijks', C, z)
    return h1, jnp.einsum('sl,ijks->ijkl', C, z)

h1_mo, eri_mo = mo_ints(C, h1e, eri)
_, ci = cs_casci(h1_mo[:nocc, :nocc],
                 eri_mo[:nocc, :nocc, :nocc, :nocc],
                 ncore=ncore, nact=nact, nele_act=nele,
                 nroots=1, return_ci=True)
C2, ci2, eps = canonicalize_jax(C, h1e, eri, ci, ncore=ncore,
                                nact=nact, nele_act=nele)
h1_mo2, eri_mo2 = mo_ints(C2, h1e, eri)
dets = cas_dets_for2(ncore, nact, nele)

print('CASSCF', mc.e_tot, 'CI norm', float(jnp.linalg.norm(ci2)))
print('eps', np.asarray(eps))
\end{lstlisting}

\texttt{dets} is a list of determinant bitmaps with the core spin
orbitals first.  \texttt{eps} is the semicanonical Fock diagonal; it
must not be replaced by the RHF orbital energies.

\subsection{The diagonal path and its class table}

\begin{lstlisting}[language=Python]
from diffmrc.core.caspt2 import caspt2_contracted_e2

e2, breakdown = caspt2_contracted_e2(
    np.asarray(h1_mo2), np.asarray(eri_mo2), np.asarray(ci2), dets,
    ncore=ncore, nact=nact, nele=nele, eps=np.asarray(eps))

for name, value in breakdown.items():
    print(f'{name:10s} {value: .10f}')
print('total', e2)
\end{lstlisting}

For H$_2$O/6-31G CAS(2,2) the two three-active-index classes
\texttt{Vm1\_a} and \texttt{V1\_i} are below $10^{-9}$~Eh: they vanish
on a two-orbital active space.  The five double and single classes and
the \texttt{V0\_ia} block carry the correction.  On CO/6-31G CAS(2,2)
the same structure holds with different weights.

\subsection{How to read the class table}

For a class $c$ with energy $e_{2,c}$ and total $e_2$, the two useful
scalars are the signed share and the absolute share,
\begin{equation}
  s_c = \frac{e_{2,c}}{e_2},
  \qquad
  a_c = \frac{|e_{2,c}|}{\sum_{c'}|e_{2,c'}|}.
  \label{eq:tutorial-class-shares}
\end{equation}
The signed share says whether the class lowers or raises the energy
relative to the total; the absolute share says which class dominates
the correlation.  A large signed share with a small absolute share is
possible when the total is small.  A cancellation index
\begin{equation}
  C = \frac{\sum_c |e_{2,c}|}{|e_2|}
  \label{eq:tutorial-class-cancellation}
\end{equation}
near one means the classes add coherently; $C\gg1$ means the total is
a cancellation between classes and the total alone is not a safe
diagnostic.

The point is not academic.  On N$_2$/6-31G CAS(6,6) the default-path
total is $0.989$ of the reference, but the class residuals are
\texttt{V1\_ija} $+55\%$, \texttt{V2\_ij} $-76\%$, \texttt{V0\_ia}
$+10\%$, and \texttt{Vm2\_ab} $+12\%$.  The total is close for the
wrong reason; the class table is what identifies the cancellation.
On the CAS(2,2) anchors the class ratios are all within a few percent
of one, so the total is a meaningful summary there.

The denominator sign matters too.  A class with $A-E_0S>0$ lowers the
energy; a class with $A-E_0S<0$ would raise it.  The code skips a
denominator with $|A-E_0S|<10^{-10}$; a skipped class is reported as
zero in the breakdown, not as a missing key.  When a class is
unexpectedly zero, print the denominator and the overlap eigenvalue
before attributing it to symmetry.

The default \texttt{class\_diagonal='function'} evaluates each
contracted function with its own denominator.  When the class overlap
has degenerate eigenvalues the diagonal reading depends on the basis
chosen inside the degenerate block.  The deterministic gauge makes the
number reproducible, but the class energy is invariant only under
\texttt{class\_diagonal='block'}:

\begin{lstlisting}[language=Python]
e2_block, bd_block = caspt2_contracted_e2(
    np.asarray(h1_mo2), np.asarray(eri_mo2), np.asarray(ci2), dets,
    ncore=ncore, nact=nact, nele=nele, eps=np.asarray(eps),
    class_diagonal='block', v0_ia='class')

for name in breakdown:
    print(f'{name:10s} default {breakdown[name]: .8f}'
          f'   block {bd_block[name]: .8f}')
\end{lstlisting}

On the CAS(2,2) anchors the block and per-function readings lie close
together.  The documented machine-precision agreement is between the
per-function reading and the orthonormalized diagonal
(\texttt{class\_diagonal='ortho'}); on larger active spaces the
readings separate, and the N$_2$ case in the manual chapter is the one
that separates them.

\subsection{The coupled full-Fock path}

\begin{lstlisting}[language=Python]
from diffmrc.core.caspt2 import caspt2_contracted_e2_coupled

e2_cpl, bd_cpl, system = caspt2_contracted_e2_coupled(
    np.asarray(h1_mo2), np.asarray(eri_mo2), np.asarray(ci2), dets,
    ncore=ncore, nact=nact, nele=nele, return_system=True)

print('coupled', e2_cpl)
print('labels ', system['labels'])
print('E0     ', system['E0'])
print('x norm ', float(np.linalg.norm(system['x'])))
print('min eig(S)', float(np.min(np.linalg.eigvalsh(system['S']))))
\end{lstlisting}

The coupled path keeps the inter-class coupling of the state-specific
generalized Fock operator.  It is a NumPy value layer; it does not
replace the differentiable diagonal path.  The \texttt{system}
dictionary is the object to inspect when a class value is questioned:
the overlap spectrum, the reduced solution \texttt{x}, and the class
labels are all there.

\subsection{The three-center path}

The three-center contraction consumes the metric-folded factor
\texttt{u\_mo} instead of the four-center tensor:

\begin{lstlisting}[language=Python]
from diffmrc.core import ri
from diffmrc.core.caspt2_three_center import (
    caspt2_contracted_e2_three_center)

auxmol, label = ri.make_auxmol(mol)
factors = ri.ri_factors(mol, auxmol)
u_mo, _ = ri.ri_eri_mo(np.asarray(C2), factors.b, factors.jhalf,
                       full=False)

e2_3c, bd_3c = caspt2_contracted_e2_three_center(
    np.asarray(h1_mo2), u_mo, np.asarray(ci2), dets,
    ncore=ncore, nact=nact, nele=nele, eps=np.asarray(eps))

print(label)
print('three-center', e2_3c)
print('max class difference',
      max(abs(bd_3c[k] - breakdown[k]) for k in breakdown))
\end{lstlisting}

The two paths use the same class definitions and the same relative
weight between the two orderings, so the class difference is at the
roundoff level.  The three-center path is not a different
approximation: it computes the same CASPT2-D energy with the RI
factorization of the integrals.  The RI error enters through
\texttt{u\_mo}, as it does in the RI-CASCI chapter.

\subsection{The N\texorpdfstring{$_2$}{2} case}

The system that exercises the three-active-index classes is
N$_2$/6-31G CAS(6,6).  It needs a converged CASSCF with the NR orbital
step and a tight gradient tolerance; the reference values used by the
validation are RHF $-108.867763375101$, CASSCF
$-109.01554685351947$, CASPT2-D $-0.071520733$, and full-FIC
$-0.07301141705494$~Eh.  The driver in
\texttt{scripts/caspt2\_n2\_validation.py} splits the reference build
and the two expensive contractions into separate stages:

\begin{lstlisting}
python scripts/caspt2_n2_validation.py --stage ref --outdir out
python scripts/caspt2_n2_validation.py --stage default --outdir out
python scripts/caspt2_n2_validation.py --stage coupled --outdir out
\end{lstlisting}

The default-path ratio to the frozen ORCA CASPT2-D reference is
$0.989$; against the release reference at one BLAS thread it is
$0.983$.  The two three-active-index classes contribute $-0.003648$ and
$-0.001625$~Eh.  The coupled path reaches $1.0232$.  The residual is a
class-space difference in \texttt{V1\_ija}, documented in the manual
chapter, not a convergence artifact.

\subsection{What to check before trusting a number}

\begin{itemize}
\item The CI vector is state-specific and matches the determinant list.
\item \texttt{eps} comes from \texttt{canonicalize\_jax}, not from the
  SCF.
\item The class table is read, not just the total.  A total that looks
  close to a reference can be a cancellation between classes with
  opposite errors.
\item When a class value enters a decision, use
  \texttt{class\_diagonal='block'} and \texttt{v0\_ia='class'}; the
  default diagonal is gauge dependent on degenerate class overlaps.
\item The three-center path needs the metric-folded factor; the raw
  \texttt{ri\_three\_center} tensor is not the input.
\end{itemize}
